%% This file contains a number of tweaks that are typically applied to the main document.
%% They are not enabled by default, but can be enabled by uncommenting the relevant lines.

%%
%% Inline annotations; for predefined colors, refer to "dvipsnames" in the xcolor package:
%% https://tinyurl.com/overleaf-colors
%%
\newcommand{\red}[1]{{\color{red}#1}}
\newcommand{\todo}[1]{{\color{red}#1}}
\newcommand{\TODO}[1]{\textbf{\color{red}[TODO: #1]}}


%%% NEW PACKAGES

\usepackage{float}
\usepackage{siunitx}
\usepackage{xspace}
\usepackage{adjustbox}
\usepackage{multirow}
\usepackage{color, colortbl}
\usepackage{makecell}

\usepackage[accsupp]{axessibility}

\newcommand{\methodname}{Vista4D\xspace}
\newcommand{\website}{https://eyeline-labs.github.io/Vista4D}
\newcommand{\projectpage}{\href{\website}{project page}}

\newcommand{\Paragraphnoskip}[1]{\noindent \textbf{#1.}}
\newcommand{\Paragraph}[1]{\smallskip\noindent \textbf{#1.}}

\newcommand{\ning}[1]{\textcolor{red}{\textbf{Ning:} #1}}
\newcommand{\jordan}[1]{\textcolor{gray}{\textbf{J:} #1}}
\newcommand{\yash}[1]{\textcolor{cyan}{\textbf{Y:} #1}}
\newcommand{\zhizheng}[1]{\textcolor{teal}{\textbf{Z:} #1}}
\newcommand{\yiwei}[1]{\textcolor{orange}{\textbf{YZ:} #1}}

\newcommand{\R}{\mathbb{R}}

\newcommand{\NN}{\mathcal{N}}
\newcommand{\LL}{\mathcal{L}}

\renewcommand{\vec}{\mathbf}
\newcommand{\src}{\mathrm{src}}
\newcommand{\tgt}{\mathrm{tgt}}
\newcommand{\gen}{\mathrm{gen}}
\newcommand{\rdr}{\mathrm{rdr}}
\newcommand{\wld}{\mathrm{wld}}
\newcommand{\dyn}{\mathrm{dyn}}
\newcommand{\stc}{\mathrm{stc}}
\newcommand{\img}{\mathrm{img}}
\newcommand{\wlddyn}{\mathrm{wld, dyn}}
\newcommand{\wldstc}{\mathrm{wld, stc}}
\newcommand{\wldtmp}{\mathrm{wld, tmp}}
\newcommand{\srctgt}{\mathrm{src \to tgt}}
\newcommand{\srctgtsrc}{\mathrm{src \to tgt \to src}}
\newcommand{\tgtsrc}{\mathrm{tgt \to src}}
\newcommand{\tgtsrctgt}{\mathrm{tgt \to src \to tgt}}

\newcommand{\tss}{\textsuperscript}
\newcommand{\specialfootnotetext}[1]{
  \begingroup
  \def\thefootnote{\fnsymbol{footnote}}  % Use symbol for footnote label
  \footnotetext[0]{#1}  % Use 0 as manual numbering, won't appear in main body
  \endgroup
}

\newcommand{\up}{$\Uparrow$}
\newcommand{\down}{$\Downarrow$}

\newcommand{\textbi}[1]{\textbf{\textit{#1}}}

\setlength{\textfloatsep}{10pt plus 1pt minus 2pt}
\setlength{\dbltextfloatsep}{10pt plus 2pt minus 2pt}
% \setlength{\floatsep}{10pt plus 1pt minus 2pt}
% \setlength{\dblfloatsep}{10pt plus 2pt minus 2pt}
\setlength{\abovecaptionskip}{7pt plus 1pt minus 2pt}

%%
%% work harder in optimizing text layout. Typically shrinks text by 1/6 of page, enable
%% it at the very end of the writing process, when you are just above the page limit
%%
\usepackage{microtype}

%%
%% fine-tune paragraph spacing
%%
% \renewcommand{\paragraph}[1]{\vspace{.5em}\noindent\textbf{#1.}}

%%
%% globally adjusts space between figure and caption
%%
% \setlength{\abovecaptionskip}{.5em}


%%
%% Allows "the use of \paper to refer to the project name"
%% with automatic management of space at the end of the word
%%
% \usepackage{xspace}
% \newcommand{\paper}{ProjectName\xspace}

%%
%% Commonly used math definitions
%%
% \DeclareMathOperator*{\argmin}{arg\,min}
% \DeclareMathOperator*{\argmax}{arg\,max}

%%
%% Tigthen underline
%%
% \usepackage{soul}
% \setuldepth{foobar}